\chapter[Introducción a Bases de Datos]{\emj{📚}~Introducción a Bases de Datos}

\section[¿Qué es una base de datos?]{\emj{🧠}~¿Qué es una base de datos?}

Una base de datos es un sistema organizado para almacenar, gestionar y recuperar información de forma persistente y estructurada.
La clave es la palabra \emph{organizado}: los datos no solo se guardan, sino que se almacenan con una estructura que permite recuperarlos eficientemente, garantiza su integridad y permite consultarlos de formas arbitrarias.

El \textbf{DBMS} (Database Management System) es el software que se interpone entre la aplicación y el almacenamiento físico.
La aplicación nunca toca directamente los archivos en disco; todo pasa por el DBMS, que administra acceso concurrente, integridad de datos, recuperación ante fallos y seguridad.
Esta separación de responsabilidades es fundamental: permite cambiar el motor de base de datos sin modificar la lógica de negocio de la aplicación.

\begin{teoriabox}{Componentes internos de un DBMS}
\begin{itemize}
  \item \textbf{Motor de almacenamiento:} decide cómo se organizan los datos en disco (páginas, extents, tablespaces). InnoDB y MyISAM son motores distintos en MySQL.
  \item \textbf{Gestor de consultas (Query Planner):} analiza el SQL, genera un árbol de operaciones y elige el plan de ejecución más barato según estadísticas.
  \item \textbf{Gestor de transacciones:} coordina el acceso concurrente garantizando ACID. Usa locks o MVCC para aislar transacciones.
  \item \textbf{Gestor de buffer (Buffer Pool):} mantiene páginas de disco en memoria RAM para evitar lecturas físicas costosas. Es el componente más crítico para el rendimiento.
  \item \textbf{Gestor de recuperación (WAL):} escribe un log de cambios antes de modificar los datos. Si el servidor se cae, el log permite reconstruir el estado consistente.
  \item \textbf{BD:} colección estructurada de datos relacionados con semántica definida por el esquema
  \item \textbf{RDBMS:} DBMS basado en modelo relacional — PostgreSQL, MySQL, Oracle, SQL Server
\end{itemize}
\end{teoriabox}

\section[Tipos de bases de datos]{\emj{📋}~Tipos de bases de datos}

No existe un tipo de base de datos universalmente superior. Cada modelo está optimizado para un patrón de acceso específico y un tipo de dato diferente.
Elegir el modelo equivocado puede resultar en un sistema que funciona pero que es imposible de escalar o mantener a mediano plazo.
La tendencia moderna es \textbf{polyglot persistence}: usar múltiples tipos de BD en el mismo sistema, cada uno para lo que hace mejor.

\begin{itemize}
  \item \textbf{Relacional (SQL):} tablas con esquema fijo, joins entre tablas, transacciones ACID, lenguaje SQL estándar. Ideal para datos estructurados con relaciones complejas. PostgreSQL, MySQL, Oracle, SQL Server.
  \item \textbf{Documental:} documentos JSON/BSON, esquema flexible por documento, bueno para datos heterogéneos o con estructura variable. MongoDB, CouchDB, Firestore.
  \item \textbf{Clave-valor:} acceso O(1) por clave, sin estructura interna de valor, extremadamente rápido en lectura. Redis, DynamoDB, Memcached.
  \item \textbf{Columnar:} optimizado para leer muchas filas de pocas columnas (analítica), no para transacciones. Cassandra, ClickHouse, Apache HBase.
  \item \textbf{Grafos:} entidades como nodos, relaciones como aristas con propiedades. Ideal cuando las relaciones son tan importantes como los datos. Neo4j, Amazon Neptune.
  \item \textbf{Series de tiempo:} datos con timestamp, optimizado para ingestión masiva y consultas por rango temporal. InfluxDB, TimescaleDB (extensión de PG).
\end{itemize}

\begin{tipbox}{¿Cuándo elegir SQL vs NoSQL?}
Usa SQL (relacional) por defecto: es maduro, tiene garantías ACID, soporte de joins y herramientas consolidadas.
Considera NoSQL cuando: el esquema varía mucho entre registros (documental), necesitas caché de baja latencia (clave-valor), escalado horizontal masivo de escritura (columnar), o las relaciones entre entidades son el dato principal (grafos).
En sistemas complejos, la respuesta es frecuentemente ``ambos''.
\end{tipbox}

\section[DBMS vs archivo plano]{\emj{⚖️}~DBMS vs archivo plano}

Es válido preguntarse por qué no usar simplemente archivos CSV o JSON en disco.
La respuesta es la \textbf{concurrencia}: cuando dos procesos leen y escriben un archivo simultáneamente sin coordinación, los datos se corrompen.
Un DBMS resuelve este problema y muchos más que no son obvios hasta que ocurren en producción.

Imagine un sistema de inventario con archivos CSV: si dos cajeros descuentan el último artículo al mismo tiempo, ambos leen ``stock: 1'', ambos calculan ``stock: 0'' y ambos escriben. El resultado es stock negativo.
Un DBMS con transacciones garantiza que solo uno de los dos vea el artículo disponible.

\begin{itemize}
  \item \textbf{Concurrencia controlada:} múltiples usuarios simultáneos sin corrupción de datos (locks, MVCC)
  \item \textbf{Integridad referencial:} las FK garantizan que no existan huérfanos (pedidos sin cliente)
  \item \textbf{Recuperación ante fallos:} WAL/journaling reconstruye el estado consistente tras una caída
  \item \textbf{Consultas declarativas:} SQL describe \emph{qué} quieres; el motor elige \emph{cómo} obtenerlo
  \item \textbf{Seguridad granular:} permisos a nivel de tabla, columna o fila por usuario/rol
  \item \textbf{Índices automáticos:} acceso eficiente sin que la aplicación gestione estructuras de búsqueda
\end{itemize}

\begin{lstlisting}[caption={DDL + DML básico}]
-- Crear tabla
CREATE TABLE empleados (
    id        SERIAL PRIMARY KEY,
    nombre    VARCHAR(100) NOT NULL,
    depto     VARCHAR(50),
    salario   NUMERIC(10,2) DEFAULT 0
);

-- Insertar
INSERT INTO empleados (nombre, depto, salario)
VALUES ('Ana López', 'Ingeniería', 45000.00);

-- Consultar
SELECT nombre, salario
FROM   empleados
WHERE  depto = 'Ingeniería'
ORDER BY salario DESC;
\end{lstlisting}

\begin{lstlisting}[caption={Tipos de datos comunes}]
-- Numéricos
id      SERIAL          -- entero autoincremental (alias de INTEGER + secuencia)
precio  NUMERIC(10,2)   -- exacto: 10 dígitos, 2 decimales
ratio   FLOAT           -- punto flotante (impreciso para dinero)

-- Texto
nombre  VARCHAR(100)    -- longitud máxima 100
codigo  CHAR(6)         -- longitud fija 6 (rellena con espacios)
notas   TEXT            -- longitud ilimitada

-- Fechas
creado  TIMESTAMP DEFAULT NOW()
fecha   DATE
\end{lstlisting}

\begin{warningbox}{FLOAT para dinero}
Nunca uses \texttt{FLOAT} o \texttt{DOUBLE} para valores monetarios. El estándar IEEE 754 de punto flotante introduce errores de representación binaria: \texttt{0.1 + 0.2 = 0.30000000000000004}.
En operaciones financieras esto acumula errores que producen discrepancias al centavo.
Usa siempre \texttt{NUMERIC(p,s)} o \texttt{DECIMAL(p,s)}, que almacena el número exacto en base 10.
\end{warningbox}

\begin{entrevistabox}
\begin{enumerate}
  \item ¿Cuál es la diferencia entre DBMS y RDBMS?
  \item ¿Cuándo elegirías NoSQL sobre SQL? Da un ejemplo concreto.
  \item ¿Qué ventajas ofrece un DBMS frente a archivos planos para una aplicación con múltiples usuarios?
\end{enumerate}
\end{entrevistabox}
